\begin{table}

\caption{Party Supply by Country for Pro-Normalization Parties. Average position indicates the average score on the defence normalization index of respondents in that country. The number of pro-normalization parties indicates that number of parties that score above a 5 from favouring close ties with Russia in the 2023 CHES dataset. Right-wing parties are those that are more TAN from the CHES dataset. Polling data is drawn from Politico's Poll of Polls in 2025.\label{tab:party_supply}}
\centering
\begin{tabular}[t]{lcclcc}
\toprule
Country & Average Position & \textbackslash{}\# Normalization Parties & Ideology & Polling \textbackslash{}\% 2025 & Spread of Pro-Russia Party Scores\\
\midrule
Bulgaria & 0.78 & 2 & RW:2; LW:0 & 34.02\% & 3.65\\
Greece & 0.64 & 4 & RW:4; LW:0 & 33.25\% & 3.18\\
Hungary & 0.63 & 2 & RW:2; LW:0 & 42.06\% & 2.87\\
Slovakia & 0.47 & 4 & RW:4; LW:0 & 33.53\% & 3.67\\
Italy & 0.46 & 3 & RW:2; LW:1 & 49.39\% & 2.26\\
Romania & 0.41 & 1 & RW:1; LW:0 & 46.56\% & 1.61\\
Croatia & 0.19 & 1 & RW:1; LW:0 & No Parties & 1.66\\
France & 0.11 & 4 & RW:3; LW:1 & 34.44\% & 1.95\\
Germany & 0.00 & 2 & RW:1; LW:1 & 38.59\% & 3.74\\
Spain & -0.10 & 4 & RW:1; LW:3 & 12.78\% & 1.60\\
Poland & -0.22 & 1 & RW:1; LW:0 & No Parties & 2.66\\
Netherlands & -0.26 & 3 & RW:2; LW:1 & 21.99\% & 2.51\\
UK & -0.27 & 0 & No Parties & No Parties & 0.47\\
Denmark & -0.39 & 0 & No Parties & No Parties & 0.36\\
Lithuania & -0.43 & 2 & RW:2; LW:0 & 20.95\% & 3.41\\
Sweden & -0.44 & 0 & No Parties & 24.59\% & 1.13\\
Finland & -0.59 & 0 & No Parties & No Parties & 0.56\\
\bottomrule
\end{tabular}
\end{table}
